% Preamble: \usepackage{circuitikz,siunitx,subcaption}
% Include from the repository root with % Preamble: \usepackage{circuitikz,siunitx,subcaption}
% Include from the repository root with % Preamble: \usepackage{circuitikz,siunitx,subcaption}
% Include from the repository root with % Preamble: \usepackage{circuitikz,siunitx,subcaption}
% Include from the repository root with \input{aula6/circuitos}.
\begin{figure}[htbp]
    \centering
    \captionsetup[subfigure]{labelformat=simple, labelsep=space}
    \renewcommand{\thesubfigure}{\alph{subfigure})}

    \begin{minipage}{0.8\linewidth}
    \centering
    \begin{subfigure}[t]{0.48\linewidth}
        \centering
        \resizebox{\linewidth}{!}{\input{aula6/retificador_meia_onda}}
        \caption{Retificador de meia onda.}
        \label{circ:retificador-meia-onda}
    \end{subfigure}\hfill%
    \begin{subfigure}[t]{0.48\linewidth}
        \centering
        \resizebox{\linewidth}{!}{\input{aula6/retificador_onda_completa}}
        \caption{Retificador de onda completa.}
        \label{circ:retificador-onda-completa}
    \end{subfigure}

    \bigskip
    \begin{subfigure}[t]{0.48\linewidth}
        \centering
        \resizebox{\linewidth}{!}{\input{aula6/ceifador_serie}}
        \caption{Ceifador em série.}
        \label{circ:ceifador-serie}
    \end{subfigure}\hfill%
    \begin{subfigure}[t]{0.48\linewidth}
        \centering
        \resizebox{\linewidth}{!}{\input{aula6/ceifador_paralelo}}
        \caption{Ceifador em paralelo.}
        \label{circ:ceifador-paralelo}
    \end{subfigure}
    \end{minipage}

    \caption{Circuitos retificadores e ceifadores.}
    \label{fig:circuitos-aula6}
\end{figure}
.
\begin{figure}[htbp]
    \centering
    \captionsetup[subfigure]{labelformat=simple, labelsep=space}
    \renewcommand{\thesubfigure}{\alph{subfigure})}

    \begin{minipage}{0.8\linewidth}
    \centering
    \begin{subfigure}[t]{0.48\linewidth}
        \centering
        \resizebox{\linewidth}{!}{\begin{tikzpicture}
    % Shared bounds keep all four drawings at the same scale.
    \path[use as bounding box] (-1.6, -0.5) rectangle (7.6, 4.8);
    \draw (0, 0) to[sinusoidal voltage source,
        l={\shortstack{$20\ \si{\volt}$\\$60\ \si{\hertz}$}},
        v_>={$V_{\mathrm{in}}$},
        voltage/distance from node=0.3, voltage/bump b=2.2] (0, 4);

    \draw (0, 4) to[empty diode, l={$\mathrm{IN4148}$}] (6, 4);
    \draw (6, 4) to[american resistor, l_={$10\ \si{\kilo\ohm}$},
        v^<={$V_{\mathrm{out}}$},
        voltage/distance from node=0.4, voltage/bump b=2.1] (6, 0);
    \draw (6, 0) -- (0, 0);
\end{tikzpicture}
}
        \caption{Retificador de meia onda.}
        \label{circ:retificador-meia-onda}
    \end{subfigure}\hfill%
    \begin{subfigure}[t]{0.48\linewidth}
        \centering
        \resizebox{\linewidth}{!}{\begin{tikzpicture}
    % Shared bounds keep all four drawings at the same scale.
    \path[use as bounding box] (-1.6, -0.5) rectangle (7.6, 4.8);
    \draw (0, 0) to[sinusoidal voltage source,
        l={\shortstack{$20\ \si{\volt}$\\$60\ \si{\hertz}$}},
        v_>={$V_{\mathrm{in}}$},
        voltage/distance from node=0.3, voltage/bump b=2.2] (0, 4);

    \draw (0, 4) -- (6, 4);
    \draw (0, 0) -- (6, 0);
    \draw (2.5, 2) to[empty diode, l_={$\mathrm{IN4148}$}] (2.5, 4);
    \draw (6, 4) to[empty diode, l={$\mathrm{IN4148}$}] (6, 2);
    \draw (2.5, 2) to[empty diode, l_={$\mathrm{IN4148}$}] (2.5, 0);
    \draw (6, 0) to[empty diode, l_={$\mathrm{IN4148}$}] (6, 2);
    \draw (6, 2) to[american resistor, l_={$10\ \si{\kilo\ohm}$},
        v^<={$V_{\mathrm{out}}$},
        voltage/distance from node=0.4, voltage/bump b=2.1] (2.5, 2);

    \node[circ] at (2.5, 4) {};
    \node[circ] at (2.5, 0) {};
    \node[circ] at (2.5, 2) {};
    \node[circ] at (6, 2) {};
\end{tikzpicture}
}
        \caption{Retificador de onda completa.}
        \label{circ:retificador-onda-completa}
    \end{subfigure}

    \bigskip
    \begin{subfigure}[t]{0.48\linewidth}
        \centering
        \resizebox{\linewidth}{!}{\begin{tikzpicture}
    % Shared bounds keep all four drawings at the same scale.
    \path[use as bounding box] (-1.6, -0.5) rectangle (7.6, 4.8);
    \draw (0, 0) to[sinusoidal voltage source,
        l={\shortstack{$20\ \si{\volt}$\\$60\ \si{\hertz}$}},
        v_>={$V_{\mathrm{in}}$},
        voltage/distance from node=0.3, voltage/bump b=2.2] (0, 4);

    \draw (0, 4) to[american voltage source, l={$V_{\mathrm{CC}}$}] (3, 4);
    \draw (3, 4) to[empty diode, l={$\mathrm{IN4148}$}] (6, 4);
    \draw (6, 4) to[american resistor, l_={$10\ \si{\kilo\ohm}$},
        v^<={$V_{\mathrm{out}}$},
        voltage/distance from node=0.4, voltage/bump b=2.1] (6, 0);
    \draw (6, 0) -- (0, 0);
\end{tikzpicture}
}
        \caption{Ceifador em série.}
        \label{circ:ceifador-serie}
    \end{subfigure}\hfill%
    \begin{subfigure}[t]{0.48\linewidth}
        \centering
        \resizebox{\linewidth}{!}{\begin{tikzpicture}
    % Shared bounds keep all four drawings at the same scale.
    \path[use as bounding box] (-1.6, -0.5) rectangle (7.6, 4.8);
    \draw (0, 0) to[sinusoidal voltage source,
        l={\shortstack{$20\ \si{\volt}$\\$60\ \si{\hertz}$}},
        v_>={$V_{\mathrm{in}}$},
        voltage/distance from node=0.3, voltage/bump b=2.2] (0, 4);

    \draw (0, 4) to[american resistor, l={$10\ \si{\kilo\ohm}$}] (6, 4);
    \draw (6, 2) to[empty diode, l={$\mathrm{IN4148}$}] (6, 4);
    \draw (6, 2) to[american voltage source, l_={$V_{\mathrm{CC}}$}] (6, 0);
    % Offset the voltage arrow from the wire and inset both ends from the nodes.
    \draw[-latex] (6.5, 0.4) .. controls (7, 1.2) and (7, 2.8) .. (6.5, 3.6)
        node[midway, right] {$V_{\mathrm{out}}$};
    \draw (6, 0) -- (0, 0);
\end{tikzpicture}
}
        \caption{Ceifador em paralelo.}
        \label{circ:ceifador-paralelo}
    \end{subfigure}
    \end{minipage}

    \caption{Circuitos retificadores e ceifadores.}
    \label{fig:circuitos-aula6}
\end{figure}
.
\begin{figure}[htbp]
    \centering
    \captionsetup[subfigure]{labelformat=simple, labelsep=space}
    \renewcommand{\thesubfigure}{\alph{subfigure})}

    \begin{minipage}{0.8\linewidth}
    \centering
    \begin{subfigure}[t]{0.48\linewidth}
        \centering
        \resizebox{\linewidth}{!}{\begin{tikzpicture}
    % Shared bounds keep all four drawings at the same scale.
    \path[use as bounding box] (-1.6, -0.5) rectangle (7.6, 4.8);
    \draw (0, 0) to[sinusoidal voltage source,
        l={\shortstack{$20\ \si{\volt}$\\$60\ \si{\hertz}$}},
        v_>={$V_{\mathrm{in}}$},
        voltage/distance from node=0.3, voltage/bump b=2.2] (0, 4);

    \draw (0, 4) to[empty diode, l={$\mathrm{IN4148}$}] (6, 4);
    \draw (6, 4) to[american resistor, l_={$10\ \si{\kilo\ohm}$},
        v^<={$V_{\mathrm{out}}$},
        voltage/distance from node=0.4, voltage/bump b=2.1] (6, 0);
    \draw (6, 0) -- (0, 0);
\end{tikzpicture}
}
        \caption{Retificador de meia onda.}
        \label{circ:retificador-meia-onda}
    \end{subfigure}\hfill%
    \begin{subfigure}[t]{0.48\linewidth}
        \centering
        \resizebox{\linewidth}{!}{\begin{tikzpicture}
    % Shared bounds keep all four drawings at the same scale.
    \path[use as bounding box] (-1.6, -0.5) rectangle (7.6, 4.8);
    \draw (0, 0) to[sinusoidal voltage source,
        l={\shortstack{$20\ \si{\volt}$\\$60\ \si{\hertz}$}},
        v_>={$V_{\mathrm{in}}$},
        voltage/distance from node=0.3, voltage/bump b=2.2] (0, 4);

    \draw (0, 4) -- (6, 4);
    \draw (0, 0) -- (6, 0);
    \draw (2.5, 2) to[empty diode, l_={$\mathrm{IN4148}$}] (2.5, 4);
    \draw (6, 4) to[empty diode, l={$\mathrm{IN4148}$}] (6, 2);
    \draw (2.5, 2) to[empty diode, l_={$\mathrm{IN4148}$}] (2.5, 0);
    \draw (6, 0) to[empty diode, l_={$\mathrm{IN4148}$}] (6, 2);
    \draw (6, 2) to[american resistor, l_={$10\ \si{\kilo\ohm}$},
        v^<={$V_{\mathrm{out}}$},
        voltage/distance from node=0.4, voltage/bump b=2.1] (2.5, 2);

    \node[circ] at (2.5, 4) {};
    \node[circ] at (2.5, 0) {};
    \node[circ] at (2.5, 2) {};
    \node[circ] at (6, 2) {};
\end{tikzpicture}
}
        \caption{Retificador de onda completa.}
        \label{circ:retificador-onda-completa}
    \end{subfigure}

    \bigskip
    \begin{subfigure}[t]{0.48\linewidth}
        \centering
        \resizebox{\linewidth}{!}{\begin{tikzpicture}
    % Shared bounds keep all four drawings at the same scale.
    \path[use as bounding box] (-1.6, -0.5) rectangle (7.6, 4.8);
    \draw (0, 0) to[sinusoidal voltage source,
        l={\shortstack{$20\ \si{\volt}$\\$60\ \si{\hertz}$}},
        v_>={$V_{\mathrm{in}}$},
        voltage/distance from node=0.3, voltage/bump b=2.2] (0, 4);

    \draw (0, 4) to[american voltage source, l={$V_{\mathrm{CC}}$}] (3, 4);
    \draw (3, 4) to[empty diode, l={$\mathrm{IN4148}$}] (6, 4);
    \draw (6, 4) to[american resistor, l_={$10\ \si{\kilo\ohm}$},
        v^<={$V_{\mathrm{out}}$},
        voltage/distance from node=0.4, voltage/bump b=2.1] (6, 0);
    \draw (6, 0) -- (0, 0);
\end{tikzpicture}
}
        \caption{Ceifador em série.}
        \label{circ:ceifador-serie}
    \end{subfigure}\hfill%
    \begin{subfigure}[t]{0.48\linewidth}
        \centering
        \resizebox{\linewidth}{!}{\begin{tikzpicture}
    % Shared bounds keep all four drawings at the same scale.
    \path[use as bounding box] (-1.6, -0.5) rectangle (7.6, 4.8);
    \draw (0, 0) to[sinusoidal voltage source,
        l={\shortstack{$20\ \si{\volt}$\\$60\ \si{\hertz}$}},
        v_>={$V_{\mathrm{in}}$},
        voltage/distance from node=0.3, voltage/bump b=2.2] (0, 4);

    \draw (0, 4) to[american resistor, l={$10\ \si{\kilo\ohm}$}] (6, 4);
    \draw (6, 2) to[empty diode, l={$\mathrm{IN4148}$}] (6, 4);
    \draw (6, 2) to[american voltage source, l_={$V_{\mathrm{CC}}$}] (6, 0);
    % Offset the voltage arrow from the wire and inset both ends from the nodes.
    \draw[-latex] (6.5, 0.4) .. controls (7, 1.2) and (7, 2.8) .. (6.5, 3.6)
        node[midway, right] {$V_{\mathrm{out}}$};
    \draw (6, 0) -- (0, 0);
\end{tikzpicture}
}
        \caption{Ceifador em paralelo.}
        \label{circ:ceifador-paralelo}
    \end{subfigure}
    \end{minipage}

    \caption{Circuitos retificadores e ceifadores.}
    \label{fig:circuitos-aula6}
\end{figure}
.
\begin{figure}[htbp]
    \centering
    \captionsetup[subfigure]{labelformat=simple, labelsep=space}
    \renewcommand{\thesubfigure}{\alph{subfigure})}

    \begin{minipage}{0.8\linewidth}
    \centering
    \begin{subfigure}[t]{0.48\linewidth}
        \centering
        \resizebox{\linewidth}{!}{\begin{tikzpicture}
    % Shared bounds keep all four drawings at the same scale.
    \path[use as bounding box] (-1.6, -0.5) rectangle (7.6, 4.8);
    \draw (0, 0) to[sinusoidal voltage source,
        l={\shortstack{$20\ \si{\volt}$\\$60\ \si{\hertz}$}},
        v_>={$V_{\mathrm{in}}$},
        voltage/distance from node=0.3, voltage/bump b=2.2] (0, 4);

    \draw (0, 4) to[empty diode, l={$\mathrm{IN4148}$}] (6, 4);
    \draw (6, 4) to[american resistor, l_={$10\ \si{\kilo\ohm}$},
        v^<={$V_{\mathrm{out}}$},
        voltage/distance from node=0.4, voltage/bump b=2.1] (6, 0);
    \draw (6, 0) -- (0, 0);
\end{tikzpicture}
}
        \caption{Retificador de meia onda.}
        \label{circ:retificador-meia-onda}
    \end{subfigure}\hfill%
    \begin{subfigure}[t]{0.48\linewidth}
        \centering
        \resizebox{\linewidth}{!}{\begin{tikzpicture}
    % Shared bounds keep all four drawings at the same scale.
    \path[use as bounding box] (-1.6, -0.5) rectangle (7.6, 4.8);
    \draw (0, 0) to[sinusoidal voltage source,
        l={\shortstack{$20\ \si{\volt}$\\$60\ \si{\hertz}$}},
        v_>={$V_{\mathrm{in}}$},
        voltage/distance from node=0.3, voltage/bump b=2.2] (0, 4);

    \draw (0, 4) -- (6, 4);
    \draw (0, 0) -- (6, 0);
    \draw (2.5, 2) to[empty diode, l_={$\mathrm{IN4148}$}] (2.5, 4);
    \draw (6, 4) to[empty diode, l={$\mathrm{IN4148}$}] (6, 2);
    \draw (2.5, 2) to[empty diode, l_={$\mathrm{IN4148}$}] (2.5, 0);
    \draw (6, 0) to[empty diode, l_={$\mathrm{IN4148}$}] (6, 2);
    \draw (6, 2) to[american resistor, l_={$10\ \si{\kilo\ohm}$},
        v^<={$V_{\mathrm{out}}$},
        voltage/distance from node=0.4, voltage/bump b=2.1] (2.5, 2);

    \node[circ] at (2.5, 4) {};
    \node[circ] at (2.5, 0) {};
    \node[circ] at (2.5, 2) {};
    \node[circ] at (6, 2) {};
\end{tikzpicture}
}
        \caption{Retificador de onda completa.}
        \label{circ:retificador-onda-completa}
    \end{subfigure}

    \bigskip
    \begin{subfigure}[t]{0.48\linewidth}
        \centering
        \resizebox{\linewidth}{!}{\begin{tikzpicture}
    % Shared bounds keep all four drawings at the same scale.
    \path[use as bounding box] (-1.6, -0.5) rectangle (7.6, 4.8);
    \draw (0, 0) to[sinusoidal voltage source,
        l={\shortstack{$20\ \si{\volt}$\\$60\ \si{\hertz}$}},
        v_>={$V_{\mathrm{in}}$},
        voltage/distance from node=0.3, voltage/bump b=2.2] (0, 4);

    \draw (0, 4) to[american voltage source, l={$V_{\mathrm{CC}}$}] (3, 4);
    \draw (3, 4) to[empty diode, l={$\mathrm{IN4148}$}] (6, 4);
    \draw (6, 4) to[american resistor, l_={$10\ \si{\kilo\ohm}$},
        v^<={$V_{\mathrm{out}}$},
        voltage/distance from node=0.4, voltage/bump b=2.1] (6, 0);
    \draw (6, 0) -- (0, 0);
\end{tikzpicture}
}
        \caption{Ceifador em série.}
        \label{circ:ceifador-serie}
    \end{subfigure}\hfill%
    \begin{subfigure}[t]{0.48\linewidth}
        \centering
        \resizebox{\linewidth}{!}{\begin{tikzpicture}
    % Shared bounds keep all four drawings at the same scale.
    \path[use as bounding box] (-1.6, -0.5) rectangle (7.6, 4.8);
    \draw (0, 0) to[sinusoidal voltage source,
        l={\shortstack{$20\ \si{\volt}$\\$60\ \si{\hertz}$}},
        v_>={$V_{\mathrm{in}}$},
        voltage/distance from node=0.3, voltage/bump b=2.2] (0, 4);

    \draw (0, 4) to[american resistor, l={$10\ \si{\kilo\ohm}$}] (6, 4);
    \draw (6, 2) to[empty diode, l={$\mathrm{IN4148}$}] (6, 4);
    \draw (6, 2) to[american voltage source, l_={$V_{\mathrm{CC}}$}] (6, 0);
    % Offset the voltage arrow from the wire and inset both ends from the nodes.
    \draw[-latex] (6.5, 0.4) .. controls (7, 1.2) and (7, 2.8) .. (6.5, 3.6)
        node[midway, right] {$V_{\mathrm{out}}$};
    \draw (6, 0) -- (0, 0);
\end{tikzpicture}
}
        \caption{Ceifador em paralelo.}
        \label{circ:ceifador-paralelo}
    \end{subfigure}
    \end{minipage}

    \caption{Circuitos retificadores e ceifadores.}
    \label{fig:circuitos-aula6}
\end{figure}
